\section{Geometry and Mesh}

%%%%%%%%%%%%%%%%%%%%%%%%%%%%%%%%%%%%%%%%%%%%%%%%%%%%%%%%%%%%%%%%%%%%%%%%%%%%%%%%

\subsection{Surface Curvature}

%%%%%%%%%%%%%%%%%%%%%%%%%%%%%%%%%%%%%%%%%%%%%%%%%%%%%%%%%%%%%%%%%%%%%%%%%%%%%%%%

When constructing a mesh for use with GFR, it is highly recommended to
include the high-order surface curvature information.
%
To export a high-order mesh, your mesh generation software may require
that every boundary surface be defined by a CAD or database surface, even those
that are planar surfaces.
%
Therefore, it is important to understand the requirements of your meshing
software before you begin.

Without the high-order surface curvature information, GFR will be forced
to construct the elements using straight line segments.
%
This can cause faceting of the surface shape, and
with a high-order code like GFR, this can significantly degrade the solution
accuracy in the best case scenario, with a high likelihood of instabilities.
%
For purely planar or rectangular geometries, which have no curvature,
this is not a problem.

%%%%%%%%%%%%%%%%%%%%%%%%%%%%%%%%%%%%%%%%%%%%%%%%%%%%%%%%%%%%%%%%%%%%%%%%%%%%%%%%

\subsection{Element Types}
\label{SEC:Element_Types}

%%%%%%%%%%%%%%%%%%%%%%%%%%%%%%%%%%%%%%%%%%%%%%%%%%%%%%%%%%%%%%%%%%%%%%%%%%%%%%%%

The standard element types supported by GFR are shown in
\fref{FIG:Element_Types}.
%
These include:
1D line,
2D triangle,
2D quadrilateral,
3D tetrahedron,
3D pyramid,
3D prism, and
3D hexahedron.
%
The number of \key{degrees-of-freedom (DoF)}[degree-of-freedom (DoF)],
\thatis solution points, within a given
element type is a function of the order \PO{} of the solution polynomial used.
%
\Tref{TAB:Element_DoF} lists this function for each element type and gives
example DoF counts for \p[1]--\p[4].
%
The last column of \Tref{TAB:Element_DoF} also provides the maximum solution
order that can be used for each element type.

%%%%%%%%%%%%%%%%%%%%%%%%%%%%%%%%%%%%%%%%%%%%%%%%%%%%%%%%%%%%%%%%%%%%%%%%%%%%%%%%

\begin{figure}[htb]
  \begin{centering}
    \begin{subfigure}{1.5in}
      \insertimage[width=\textwidth]{Elements/Elem_1D_Line}
      \caption{1D Line}
      \label{FIG:Elem_1D_Line}
    \end{subfigure}%
    \begin{subfigure}{1.5in}
      \insertimage[width=\textwidth]{Elements/Elem_2D_TriRt}
      \caption{2D Triangle}
      \label{FIG:Elem_2D_Tri}
    \end{subfigure}%
    \begin{subfigure}{1.5in}
      \insertimage[width=\textwidth]{Elements/Elem_2D_Quad} 
      \caption{2D Quadrilateral} 
      \label{FIG:Elem_2D_Quad} 
    \end{subfigure}
    %%%%%%%%%%%%%%%%%%%%%%%%%%%%%%%%%%%%%%%%%%%%%%%%%%%%%
    \begin{subfigure}{1.5in}
      \insertimage[width=\textwidth]{Elements/Elem_3D_TetRt}
      \caption{3D Tetrahedron}
      \label{FIG:Elem_3D_Tet}
    \end{subfigure}%
    \begin{subfigure}{1.5in}
      \insertimage[width=\textwidth]{Elements/Elem_3D_PyrTall}
      \caption{3D Pyramid}
      \label{FIG:Elem_3D_Pyr}
    \end{subfigure}
    \begin{subfigure}{1.5in}
      \insertimage[width=\textwidth]{Elements/Elem_3D_PrismRt}
      \caption{3D Prism}
      \label{FIG:Elem_3D_Pri}
    \end{subfigure}%
    \begin{subfigure}{1.5in}
      \insertimage[width=\textwidth]{Elements/Elem_3D_Hex}
      \caption{3D Hexahedron}
      \label{FIG:Elem_3D_Hex}
    \end{subfigure}
    %%%%%%%%%%%%%%%%%%%%%%%%%%%%%%%%%%%%%%%%%%%%%%%%%%%%%
    \caption{Supported element types.}
    \label{FIG:Element_Types}
  \end{centering}
\end{figure}

%%%%%%%%%%%%%%%%%%%%%%%%%%%%%%%%%%%%%%%%%%%%%%%%%%%%%%%%%%%%%%%%%%%%%%%%%%%%%%%%

\begin{table}
  % The scope of these commands is limited to within this table environment.
  % Degrees of freedom of different geometries in terms of polynomial order.
  \newcommand{\dofLine}{$(\PO+1)$}
  \newcommand{\dofTria}{$(\PO+1)(\PO+2)/2$}
  \newcommand{\dofQuad}{$(\PO+1)^2$}
  \newcommand{\dofTetr}{$(\PO+1)(\PO+2)( \PO+3)/6$}
  \newcommand{\dofPyra}{$(\PO+1)(\PO+2)(2\PO+3)/6$}
  \newcommand{\dofPris}{$(\PO+1)(\PO+1)( \PO+2)/2$}
  \newcommand{\dofHexa}{$(\PO+1)^3$}
  %
  \centering
  \caption{Degrees of freedom of supported element types.}
  \label{TAB:Element_DoF}
 %\begin{center}
    \begin{tabular}{lllrrrrc}
       Dim & Element Type  & Degrees of Freedom & % P1 &   P2 &   P3 &   P4 \\
       $\PO[1]$ & $\PO[2]$ & $\PO[3]$ & $\PO[4]$ & Maximum \PO{} \\
       \hline
       1D  & Line          & \dofLine &    2 &    3 &    4 &    5 & unlimited \\
       2D  & Triangle      & \dofTria &    3 &    6 &   10 &   15 & 8 \\
       2D  & Quadrilateral & \dofQuad &    4 &    9 &   16 &   25 & unlimited \\
       3D  & Tetrahedron   & \dofTetr &    4 &   10 &   20 &   35 & 6 \\
       3D  & Pyramid       & \dofPyra &    5 &   14 &   30 &   55 & unlimited \\
       3D  & Prism         & \dofPris &    6 &   18 &   40 &   75 & 8 \\
       3D  & Hexahedron    & \dofHexa &    8 &   27 &   64 &  125 & unlimited
    \end{tabular}
 %\end{center}
\end{table}

%%%%%%%%%%%%%%%%%%%%%%%%%%%%%%%%%%%%%%%%%%%%%%%%%%%%%%%%%%%%%%%%%%%%%%%%%%%%%%%%

\IDEA{
\begin{itemize}
\item
Describe which element types are best and why.
\item
Describe mixed element parallel efficiency.
\end{itemize}
}

%%%%%%%%%%%%%%%%%%%%%%%%%%%%%%%%%%%%%%%%%%%%%%%%%%%%%%%%%%%%%%%%%%%%%%%%%%%%%%%%

\subsection{Grid Scaling}
\label{SEC:Grid_Scaling}

%%%%%%%%%%%%%%%%%%%%%%%%%%%%%%%%%%%%%%%%%%%%%%%%%%%%%%%%%%%%%%%%%%%%%%%%%%%%%%%%

GFR expects the mesh file to be in units of meters.
%
However, in many instances the mesh file may be in another set of units
or perhaps nondimensional.
%
For these situations, the code allows the user to specify a
multiplicative scaling factor that will convert the coordinates of the
existing grid into units of meters.

\NOTE
Throughout this document, we shall refer to the units of the mesh file as
  \key*[key:GridUnits]{\emph{grid units}}[grid units]
and the units after applying the grid scaling as
  \key*[key:GridScaledUnits]{\emph{grid scaled units}}[grid scaled units]
($=[grid~units]*\inptext{Grid\_Scaling\_Factor}$).
%
For example, if the grid units are inches, the input file should include the
line
\begin{InputExample}
\item[{\inptext{Grid\_Scaling\_Factor = 0.0254~~~!~convert inches to meters}}]
\end{InputExample}

\SeeAlso{
\inp[nmlGrid::GridScalingFactor]{Grid\_Scaling\_Factor}
}

%%%%%%%%%%%%%%%%%%%%%%%%%%%%%%%%%%%%%%%%%%%%%%%%%%%%%%%%%%%%%%%%%%%%%%%%%%%%%%%%

\subsection{Mesh Resolution Requirements for Large-Eddy Simulation}
\label{SEC:Grid_LES}
\key*{}[large-eddy simulation (LES)!mesh requirements]

%In order to comprehend the mesh resolution requirements for
%\key{Large-Eddy Simulation (LES)}[large-eddy simulation],
%the user must first have some understanding of the turbulence cascade process.
%
In order to comprehend the mesh resolution requirements for
large-eddy simulation (LES), one must have some understanding of the
turbulent eddies that are being simulated.
%
Even though an instantaneous turbulent velocity field often appears chaotic,
it is actually composed of a multitude of swirling three-dimensional
turbulent eddies that are all intertwined.
%
These eddies span a continuous range of varying sizes,
from the (largest) integral length scale ($l=k/\epsilon$)
to the (smallest) Kolmogorov scale ($\eta=(\nu^3/\epsilon)^{1/4}$),
where
$k$ is the turbulent kinetic energy,
$\epsilon$ is the turbulent dissipation rate ($-dk/dt$), and
$\nu$ is the kinematic viscosity.
%
The energy contained in each eddy is proportional to its size.

%%%%%%%%%%%%%%%%%%%%%%%%%%%%%%%%%%%%%%%%%%%%%%%%%%%%%%%%%%%%%%%%%%%%%%%%%%%%%%%%

At any given point in the flow, we can perform a Fourier decomposition of
the energy containing eddies to examine the spectral distribution in terms of
wavenumber, $\kappa = 2 \pi / (\text{eddy diameter})$.
%
Doing so produces the energy spectral density, $E(\kappa)$, which is the
distribution of energy as a function of wavenumber.
%
As shown in \fref{FIG:EnergySpectra}, the horizontal axis represents
(from left to right) the range of eddy sizes from the largest to the
smallest.
%
This range increases with increasing Reynolds number.
%
The vertical axis represents the energy contained in those eddies,
and the area under the curve represents the turbulent kinetic energy
at that particular point in the flow.
%
Instabilities in the largest structures generate turbulence,
the intermediate structures transfer that energy, and
the smallest structures dissipate the energy into heat.

%%%%%%%%%%%%%%%%%%%%%%%%%%%%%%%%%%%%%%%%%%%%%%%%%%%%%%%%%%%%%%%%%%%%%%%%%%%%%%%%

\begin{figure}
  \begin{centering}
%   Sample figure:  https://eaglepubs.erau.edu/introductiontoaerospaceflightvehicles/chapter/turbulent-flows/
    \insertimage[width=5.5in]{tikz/EnergySpectra}
    \caption[Energy spectra for a turbulent flow.]%
            {Energy spectra for a turbulent flow, log-log scales.}
    \label{FIG:EnergySpectra}
  \end{centering}
\end{figure}

%%%%%%%%%%%%%%%%%%%%%%%%%%%%%%%%%%%%%%%%%%%%%%%%%%%%%%%%%%%%%%%%%%%%%%%%%%%%%%%%

The idea behind
\key*{LES}[large-eddy simulation (LES)]
is that we need not simulate the full range of scales, only the largest ones
that contain the predominant share of energy
(\forexample 80 percent of the area under the curve starting from the left).
%
Thus the grid should be fine enough to fully resolve (\thatis represent)
eddies at the highest desired wavenumber, typically within the inertial
subrange or outer part of the viscous subrange.
%
The dissipation effect of the high wavenumber (small scale) turbulence that
is not properly resolved by the grid is sometimes represented using an
explicit subgrid-scale model.
%
\key*{Implicit large-eddy simulation (ILES)}%
     [large-eddy simulation (LES)!implicit LES (ILES)]
takes a different approach, using the inherent dissipation of the
numerical scheme to act as the subgrid model.
%
For boundary layer flows at high Reynolds numbers, it may not be feasible to
perform a wall-resolved LES due to the extremely fine grid resolution
required to resolve the small-scale turbulent structures.
%
In such cases, some codes offer near-wall models or wall functions
in what is referred to as wall-modeled LES.

%%%%%%%%%%%%%%%%%%%%%%%%%%%%%%%%%%%%%%%%%%%%%%%%%%%%%%%%%%%%%%%%%%%%%%%%%%%%%%%%

In isotropic turbulence, the eddies are statistically uniform with no
preferred direction, resulting in 
$\ReyAve{\ReyFlx{u}\ReyFlx{u}}=
 \ReyAve{\ReyFlx{v}\ReyFlx{v}}=
 \ReyAve{\ReyFlx{w}\ReyFlx{w}}$.
%
This situation might occur far away from boundaries.
%
%
In boundary layer flows, the turbulent eddies are distorted by the mean shear
and physical spatial limitation in the wall-normal direction,
resulting in anisotropic turbulence that can exhibit long streamwise
vortex filaments.
%
The streamwise ($\ReyAve{\ReyFlx{u}\ReyFlx{u}}$) turbulence is accentuated
and the wall-normal ($\ReyAve{\ReyFlx{v}\ReyFlx{v}}$) component is reduced.
%
%
In free shear flows, such as jets and mixing layers, the mean shear
again causes an increase in the streamwise ($\ReyAve{\ReyFlx{u}\ReyFlx{u}}$)
component and a decrease in the shear direction
($\ReyAve{\ReyFlx{v}\ReyFlx{v}}$), but not to the same extent.
 
%%%%%%%%%%%%%%%%%%%%%%%%%%%%%%%%%%%%%%%%%%%%%%%%%%%%%%%%%%%%%%%%%%%%%%%%%%%%%%%%

\Tref{TAB:WRLES_Grid_Reqs} summarizes the grid spacing requirements for
wall-resolved LES from a number of sources.
%
These are given in terms of near-wall coordinates (\thatis
$\Delta x^+ = \Delta x \, \utau / \nu$,
where $\utau$ is the friction velocity).
%
Note that the wall-normal direction uses grid stretching with an initial
wall spacing $\Delta y^+_\text{min} \leq 1$ and 3--40 points within
$y^+ < 10$.
%
The spanwise spacing $(\Delta z^+)$ is typically 10--20, and
the streamwise spacing $(\Delta x^+)$ is often double or triple that.
%
\Linecite{Breuer_2007} notes that using overly coarse values for
$\Delta x^+$ or $\Delta z^+$ can lead to unphysical boundary layer streaks
and error in the predicted skin friction.
%
Based on these sources, for traditional finite-difference or finite-volume
solvers, a good starting grid would have
$(\Delta x^+, \Delta y^+_\text{min}, \Delta z^+) = (20\text{--}30,1,15)$
with 10 points below $y^+ < 10$
(which corresponds to a geometric growth rate of 1.02).
%
For free shear flows, 50 points or more are likely needed across the
shear layer thickness, $\delta$.
 
%%%%%%%%%%%%%%%%%%%%%%%%%%%%%%%%%%%%%%%%%%%%%%%%%%%%%%%%%%%%%%%%%%%%%%%%%%%%%%%%
%
% Our discussion at GRC recommended (x+,y+,z+)=(20-30,1,15).
%
%%%%%%%%%%%%%%%%%%%%%%%%%%%%%%%%%%%%%%%%%%%%%%%%%%%%%%%%%%%%%%%%%%%%%%%%%%%%%%%%

For GFR, the volume solution points are often defined using either
Gauss or Lobatto points, which are non-uniformly distributed toward the
boundaries of each cell (see \fref{FIG:1D_Points}).
%
In addition, the number of points in each cell is determined by the
chosen polynomial order for the solution, $\PO$.
%
When building a mesh, we therefore compare the above values to the
average solution point spacing within a high-order cell (\thatis
$\overline{\Delta x^+} = \Delta x^+_\text{cell} / (\PO+1)$).
%
%When building a mesh, we therefore recommend \emph{cell spacings} of
From that, we obtain the recommended \emph{cell spacings} of
$\Delta x^+_\text{cell}$ = (10--80)(\PO+1),
$\Delta y^+_\text{cell,min}$ = 2,
$\Delta z^+_\text{cell}$ = (10--40)(\PO+1),
and a wall-normal growth rate of 1.5 for boundary layer flows.
%
For free shear flows, we recommend isotropic cells with 
$\Delta y_\text{cell} / \delta = (\PO+1) / 50$ across the
shear layer thickness, $\delta$.
 
%%%%%%%%%%%%%%%%%%%%%%%%%%%%%%%%%%%%%%%%%%%%%%%%%%%%%%%%%%%%%%%%%%%%%%%%%%%%%%%%

% Note: The comments in the table source below provide more specific
%       reference info than just citation, to make finding that info easier
%       in the future. Some values or ranges have been rounded here for
%       consistency.
\begin{table}
\begin{center}
\caption{Grid requirements for wall-resolved large-eddy simulation.}
\label{TAB:WRLES_Grid_Reqs}
\begin{tabular}{ccccl}
  \toprule
  $\Delta x^+$ & $\Delta y^+_\text{min}$ & $\Delta z^+$ &
  $N_{pts}$, $y^+ \!\! < \!\! 10$ & Reference \\
  \midrule
  50--150 &   $< 1$  & 15--40 & 20--40 & Georgiadis\cite{Georgiadis_2009} \\
  50--150 &      1   & 10--40 &  3--5  & Breuer\cite{Breuer_2007} \\
                                       % Breuer_2007, Table 5.4, p.209.
                                       % Jiang\cite{Jiang_2016} \\
                                       % Jiang_2016, Table 2.1, p.46.
                                       %   Cites Breuer(2007).
  50--100 &      1   & 10--20 &     3  & Sagaut\cite{Sagaut_2007} \\
                                       % Sagaut_2007, Table 3.4, p.107.
  50      &      1   & 15     &        & Sagaut\cite{Sagaut_2009} \\
                                       % Sagaut_2009, p.2850.
  25      &   $< 1$  & 11     & 10     & Schlatter\cite{Schlatter_2010} \\
                                       % Schlatter_2010, p.6.
  20--30  & 0.8--1.2 & 10--15 &        & Debroeyer\cite{Debroeyer_2024} \\
                                       % Debroeyer_2024, Table 1, p.A9-4.
  10--50  & 0.5--2.5 & 10--20 &        & Ameri\cite{Ameri_2016} \\
                                       % Range of values in Table 1, p.13.
  \bottomrule
\end{tabular}
\end{center}
\end{table}

%%%%%%%%%%%%%%%%%%%%%%%%%%%%%%%%%%%%%%%%%%%%%%%%%%%%%%%%%%%%%%%%%%%%%%%%%%%%%%%%


